(1) 哈密顿量可分解为三个独立的一维振子，角频率分别为 \(\omega_x=\omega,\omega_y=2\omega,\omega_z=3\omega\)。令 \(n_x,n_y,n_z\in\{0,1,2,\ldots\}\)，能谱为 \(E_{n_xn_yn_z}=\hbar\omega[(n_x+1/2)+2(n_y+1/2)+3(n_z+1/2)]=\hbar\omega(n_x+2n_y+3n_z+3)\)，所以基态能量 \(E_0=3\hbar\omega\)，基态仅一重。频率可公度导致一些不同量子数三元组拥有相同激发能，这不同于各向同性振子的旋转简并。

(2) 激发量子数 \(N=n_x+2n_y+3n_z\) 的简并度就是该丢番图方程的非负整数解数目。母函数为 \(\sum_{N\ge0}g_Nq^N=[(1-q)(1-q^2)(1-q^3)]^{-1}\)，分别来自三方向的几何级数。对 \(N=6\)，当 \(n_z=0\) 时 \(n_x+2n_y=6\) 有4组；当 \(n_z=1\) 时余下3有2组；当 \(n_z=2\) 时余下0有1组，共 \(g_6=7\)。穷举的计数不依赖于改变能量单位或额外的自旋简并度。

(3) 设 \(x=\beta\hbar\omega\) 且 \(\beta=(k_BT)^{-1}\)。各方向能量独立，故单粒子配分函数为 \(Z=e^{-3x}/[(1-e^{-x})(1-e^{-2x})(1-e^{-3x})]\)。由 \(U=-\partial_\beta\ln Z\)，得 \(U=3\hbar\omega+\sum_{j=1}^3 j\hbar\omega/(e^{jx}-1)\)。第一项是与温度无关的基态零点能，后一项才来自各方向的热激发；将零点能误算入热容，会给出非物理的额外常数贡献。

(4) 对每个方向热能 \(U_j=j\hbar\omega/(e^{jx}-1)\) 求温度导数，得到 \(C_V=k_B\sum_{j=1}^3(jx)^2 e^{jx}/(e^{jx}-1)^2\)。当 \(x\ll1\)，每一振动方向的贡献趋近 \(k_B\)，故 \(C_V\to3k_B\)，符合三维经典谐振子每个方向动能与势能共两个二次型项的能均分。低温 \(x\gg1\) 时最低激发来自 \(x\) 方向，主导 \(C_V\sim k_Bx^2e^{-x}\)，其他 \(e^{-2x},e^{-3x}\) 项更小；第三定律要求热容趋零也得到满足。
